% (User input window)

\begin{corollary}[bulk--boundary 的双重伸缩率：边界线性斜率 \texorpdfstring{$\kappa$}{kappa}、bulk 极限斜率 \texorpdfstring{$\kappa^2$}{kappa^2}]\label{cor:xi-bulk-boundary-double-scaling-holography}
\leanverified{paper\_xi\_bulk\_boundary\_double\_scaling\_holography}
在命题 \ref{prop:xi-hankel-det-scaling-covariance} 的有限原子设定下，定义 bulk 熵势
\[
S_{\mathrm{bulk}}(m):=-\log\det \widetilde{\mathsf H}_{\kappa-1},
\]
其中 \(\widetilde{\mathsf H}_{\kappa-1}\) 为伸缩后的 Hankel 主块。
则当 \(m\to\infty\) 时有渐近
\[
\boxed{\;
S_{\mathrm{bulk}}(m)=\kappa^2\log m+O(1).\;}
\]
并且若将边界缺陷熵 \(S_{\mathrm{gen}}(m)\) 取为推论 \ref{cor:xi-pick-poisson-entropy-derivative-visible-dimension} 中的对数尺度族，则在对数尺度 \(s=\log m\to\infty\) 下成立比例律
\[
\boxed{\;
S_{\mathrm{bulk}}(e^{s})=\kappa\,S_{\mathrm{gen}}(e^{s})+O(1).\;}
\]
\end{corollary}
\begin{proof}
由命题 \ref{prop:xi-hankel-det-scaling-covariance} 的协变律，
\(
\det \widetilde{\mathsf H}_{\kappa-1}=\det \mathsf H_{\kappa-1}\cdot J_{\kappa-1}(m;\delta)
\)，
故
\(
S_{\mathrm{bulk}}(m)=S_{\mathrm{bulk}}(1)-\log J_{\kappa-1}(m;\delta)
\)。
由
\[
\log J_{\kappa-1}(m;\delta)
=\kappa(\kappa-1)\log m
+(2\kappa-1)\sum_{j=1}^{\kappa}\Bigl(\log(1+\delta_j)-\log(1+m\delta_j)\Bigr)
\]
并用渐近 \(\log(1+m\delta_j)=\log m+\log\delta_j+O(m^{-1})\)，得到
\(
\log J_{\kappa-1}(m;\delta)=-\kappa^2\log m+O(1)
\)，从而
\(
S_{\mathrm{bulk}}(m)=\kappa^2\log m+O(1)
\)。
另一方面，推论 \ref{cor:xi-pick-poisson-entropy-derivative-visible-dimension} 给出
\(
S_{\mathrm{gen}}(e^s)=S_{\mathrm{gen}}(1)+\kappa s
\)。
两式合并即得
\(
S_{\mathrm{bulk}}(e^s)-\kappa S_{\mathrm{gen}}(e^s)=O(1)
\)。证毕。
\end{proof}

\begin{corollary}[bulk--boundary 双熵的二次封口与无参审计量]\label{cor:xi-bulk-boundary-quadratic-closure-audit}
\leanverified{paper\_xi\_bulk\_boundary\_quadratic\_closure\_audit}
在推论 \ref{cor:xi-bulk-boundary-double-scaling-holography} 的有限原子设定下，
若边界缺陷熵族满足严格线性律
\(
S_{\mathrm{gen}}(m)=S_{\mathrm{gen}}(1)+\kappa\log m
\)，
则当 \(m\to\infty\) 时有二次关系
\[
\boxed{\;
S_{\mathrm{bulk}}(m)
=\frac{S_{\mathrm{gen}}(m)^2}{\log m}+O(1).\;}
\]
从而量
\[
\mathcal Q(m):=S_{\mathrm{bulk}}(m)-\frac{S_{\mathrm{gen}}(m)^2}{\log m}
\]
在 \(m\to\infty\) 下保持有界。
\end{corollary}
\begin{proof}
由推论 \ref{cor:xi-bulk-boundary-double-scaling-holography} 得
\(
S_{\mathrm{bulk}}(m)=\kappa^2\log m+O(1)
\)。
另一方面，线性律给出
\(
S_{\mathrm{gen}}(m)=\kappa\log m+O(1)
\)，故
\[
\frac{S_{\mathrm{gen}}(m)^2}{\log m}
=\frac{\bigl(\kappa\log m+O(1)\bigr)^2}{\log m}
=\kappa^2\log m+O(1).
\]
两式相减即得所示二次关系与 \(\mathcal Q(m)\) 的有界性。证毕。
\end{proof}

\begin{theorem}[bulk 熵势的可积重整化常数与 bulk--boundary 异常差极限]\label{thm:xi-bulk-renormalized-constant-and-anomaly-gap}
\leanverified{paper\_xi\_bulk\_renormalized\_constant\_and\_anomaly\_gap}
在推论 \ref{cor:xi-bulk-boundary-double-scaling-holography} 的有限原子设定与记号下，极限
\[
\mathcal{C}_{\mathrm{bulk}}(\delta)
:=\lim_{m\to\infty}\Bigl(S_{\mathrm{bulk}}(m)-\kappa^2\log m\Bigr)
\]
存在，并由显式闭式给出：
\[
\boxed{\;
\mathcal{C}_{\mathrm{bulk}}(\delta)
=S_{\mathrm{bulk}}(1)-(2\kappa-1)\sum_{j=1}^{\kappa}\log\frac{1+\delta_j}{\delta_j}\; }.
\]
进一步，若边界缺陷熵族满足严格线性律
\(
S_{\mathrm{gen}}(m)=S_{\mathrm{gen}}(1)+\kappa\log m
\)，
则 bulk--boundary 的异常差极限
\[
\mathcal{A}_{\mathrm{bulk}\mid\mathrm{gen}}(\delta)
:=\lim_{m\to\infty}\Bigl(S_{\mathrm{bulk}}(m)-\kappa S_{\mathrm{gen}}(m)\Bigr)
\]
亦存在，并满足闭式
\[
\boxed{\;
\mathcal{A}_{\mathrm{bulk}\mid\mathrm{gen}}(\delta)
=S_{\mathrm{bulk}}(1)-\kappa S_{\mathrm{gen}}(1)-(2\kappa-1)\sum_{j=1}^{\kappa}\log\frac{1+\delta_j}{\delta_j}\; }.
\]
\end{theorem}
\begin{proof}
由推论 \ref{cor:xi-bulk-boundary-double-scaling-holography} 的证明中给出的
\(
S_{\mathrm{bulk}}(m)=S_{\mathrm{bulk}}(1)-\log J_{\kappa-1}(m;\delta)
\)
与
\[
\log J_{\kappa-1}(m;\delta)
=\kappa(\kappa-1)\log m
+(2\kappa-1)\sum_{j=1}^{\kappa}\Bigl(\log(1+\delta_j)-\log(1+m\delta_j)\Bigr),
\]
并用渐近
\(
\log(1+m\delta_j)=\log m+\log\delta_j+o(1)
\)
（\(m\to\infty\)），得到
\[
\log J_{\kappa-1}(m;\delta)
=-\kappa^2\log m+(2\kappa-1)\sum_{j=1}^{\kappa}\log\frac{1+\delta_j}{\delta_j}+o(1).
\]
代回并移项即得 \(\mathcal{C}_{\mathrm{bulk}}(\delta)\) 的存在与闭式。
若进一步 \(S_{\mathrm{gen}}(m)=S_{\mathrm{gen}}(1)+\kappa\log m\)，则
\[
S_{\mathrm{bulk}}(m)-\kappa S_{\mathrm{gen}}(m)
=\Bigl(S_{\mathrm{bulk}}(m)-\kappa^2\log m\Bigr)-\kappa S_{\mathrm{gen}}(1),
\]
令 \(m\to\infty\) 并代入上一段极限即可得到 \(\mathcal{A}_{\mathrm{bulk}\mid\mathrm{gen}}(\delta)\) 的存在与闭式。证毕。
\end{proof}

\begin{remark}[边界临界的二次普适指数：碰撞触发的熵奇点系数恒为 \(2\)]\label{rem:xi-collision-universal-exponent-two}
在 bulk 口径中，命题 \ref{prop:xi-depth-hankel-determinant-vandermonde-square} 给出
\(
\det \mathsf H_{\kappa-1}
=\bigl(\prod_j\widetilde w_j\bigr)\prod_{j<k}(b_k-b_j)^2
\)；
在 boundary 口径中，命题 \ref{prop:xi-pick-poisson-det-factorization} 给出
\(
\det\mathsf P=(\prod_j p_j)\prod_{j<k}\rho_{jk}^2
\)。
因此当某对参数发生碰撞 \(b_j\to b_k\) 或 \(\rho_{jk}\to 0\)（其余保持在紧域内）时，
对应熵势出现普适的对数发散
\[
S_{\mathrm{bulk}}
:=-\log\det \mathsf H_{\kappa-1}
=2\log\frac{1}{|b_k-b_j|}+O(1),
\qquad
S_{\mathrm{gen}}
:=-\log\det\mathsf P
=2\log\frac{1}{\rho_{jk}}+O(1).
\]
该系数 \(2\) 由 Gram 行列式的平方退化结构强制给出，与权重与其余点的具体配置无关。
\end{remark}
